\section{One source window for a fixed family of cyclic hosts}
\label{sec:cubic-abelian-common-window}

\paragraph{Status and input boundary.}
This manuscript proves a finite-family strengthening of the actual
source window construction.  It is not a consequence of treating
separate instances of \texttt{VWinAt} as black boxes.
It retains the original source arithmetic and automorphic inputs
listed below.  It does not deduce those inputs from hypothesis A,
construct the mixed realization $\mathcal R_{\rm cub}$, or prove
the remainder of the index bound (TC.3).
All source paths are relative to \texttt{src/} at commit
\texttt{7cbf7a0c32fca1badf05931879a37c3f8a2b03b4}.
The proof is a mathematical replay of the displayed source definitions;
no Lean execution or independent verification of the whole repository
is asserted.

\paragraph{The exact retained source package.}
Use the literal definitions \texttt{LData}, \texttt{shaInv},
\texttt{VVat}, \texttt{VVle}, \texttt{PointHV}, and
\texttt{ThmBcFree}.  The package used here consists of
\[
 \begin{gathered}
 C_K\geq1,\qquad \texttt{LemmaK}(C_K),\qquad
 \texttt{FixedIn},\qquad\texttt{TwoIn},\\
 \texttt{PointHV},\qquad\texttt{ThmBcFree}.
 \end{gathered}                                                   \tag{AW.1}
\]
In particular the third and fourth entries have their actual
quantifier order.  For every fixed abelian field datum $M$,
\texttt{TwoIn} supplies one integer $R_0(M)$ such that for
\emph{every} admissible pair $q,d$ there is a number field $L_0$ with
\[
 \texttt{TwoBase}(M,q,d,L_0),\qquad
 \texttt{ABound}(2,R_0(M),L_0).                                  \tag{AW.2}
\]
The field $L_0$ can depend on $q,d$; the rank $R_0(M)$ does not.
More explicitly, the complete containing proof
\texttt{QuadraticExtensions/IsAlgClosureRat.lean} first obtains a
rank $g$ by Ferrero--Washington over the fixed datum
\texttt{withI}(M), then proves \texttt{baseRk}/\texttt{towerTwo}
with rank $4(g+6\deg M)$ for every such $q,d$.
Its \texttt{twoIn} passes this uniform rank to the tower-bound
input.  The ramification and tower arguments imported there are
retained; no rank for a growing field is silently replaced by
the rank of a different fixed field.
For a fixed odd prime $\ell$ and fixed datum $A$ containing
$\mu_\ell$, \texttt{FixedIn} supplies an $\texttt{ABound}$
rank for a model of $A$.  These are assertions about the actual
strict Tate groups and their cyclotomic layers.
They are not inferred here from a finite class-number calculation.

The containing file \texttt{LayeredIsomorphisms/SecondBound.lean}
states these definitions.  The complete short proof
\texttt{ThmBabOfLemmaK.lean} obtains them from the source's
Poitou--Tate and local-duality inputs, Ferrero--Washington, and
\texttt{PropA.aBound\_of\_isTower}.  The complete
\texttt{HeightAndVanishing/JointWinOfCounts.lean} proof supplies
\texttt{PointHV} from the original \texttt{CountI},
\texttt{CountII}, \texttt{CountIII}, \texttt{CountV},
\texttt{CountZ}, \texttt{LemmaK}, and duality inputs.
The \texttt{HVCore.of\_four\_alt} proof also supplies the
retained \texttt{ThmBcFree}.  The larger actual wrapper
\texttt{VwinOfRow.lean} passes the original \texttt{Row0b}
input through \texttt{automorphic\_window\_of\_inputs}, with
the stated open-image and base-change inputs.
These named upstream inputs remain the original inputs; their
deep proofs are not replaced by the argument below.

Let $S_\ell(E,e)$ denote the strict group with coefficient
$\mathbb Q_\ell/\mathbb Z_\ell(-e)$, unramified away from
$\ell$ and zero at every completion above $\ell$.
The source comparison \texttt{shaLen\_eq\_shaInv} and
\texttt{vvle\_iff\_shaInv}, together with AC-G31's literal
\texttt{LData} identification, uses precisely these groups.
Thus \texttt{VVle} bounds the nonnegative all-prime sum
\[
 \mathcal B(E,t)=
  \sum_\ell\sum_{1\leq e<t}\log\#S_\ell(E,e).                    \tag{AW.3}
\]
Equivalently one may first use every finite set of coefficient
primes, as the source does.  Finiteness and finite prime support
for fixed $E,t$ are supplied by AC-K10.

\begin{theorem}[AC-G172: a common arithmetic window for the abelian sector]
\label{AC-G172}
Retain (AW.1).  Fix a finite family $\Psi$ of primitive Dirichlet
characters, a real $C_t\geq1$, and
$0<\varepsilon<1/256$.  Let $h(u)$ be any page satisfying
(TC.2).  There are $J\geq1$, $X_0$, and $C\geq1$, independent of
$X,q,d,u$, such that for every $X\geq X_0$ and every prescribed
$\operatorname{WinPair}(\varepsilon,X^{1/J},q,d)$ there is
one $W$ satisfying
\[
 X^{1/J}\leq W,\qquad W^2\leq X,                                \tag{AW.4}
\]
and, simultaneously for every $\psi\in\Psi$ and
$u\in[W,W^2]$, the actual field
\[
 E_{\psi,u}=
 \mathbb Q(\zeta_{2^{h(u)+2}})^+
 \mathbb Q(\sqrt{-q})\mathbb Q(\sqrt{-d})L_\psi
\]
satisfies
\[
 \mathcal B(E_{\psi,u},\lfloor u^{C_t}\rfloor)
       \leq C\exp(Cu^{4\varepsilon/5})
       =o\bigl(\exp(u^\varepsilon/2)\bigr).                     \tag{AW.5}
\]
The little-o assertion is uniform on these windows and prescribed
pairs.  The constants may depend on the fixed family, retained
source inputs, $C_t$, and $\varepsilon$.

In particular take the $24$ characters of modulus $84$, represented
by their associated primitive characters, and $C_t=401$.
On the same $W$, the actual abelian field
\[
 F_u=\mathbb Q(\zeta_{2^{h(u)+2}},\zeta_7,
                   \sqrt{-q},\sqrt{-d},\omega)
\]
satisfies
\[
 \mathcal B(F_u,\lfloor u^{401}\rfloor)
               =o\bigl(\exp(u^\varepsilon/2)\bigr).             \tag{AW.6}
\]
For the original coefficient rings $O_\ell$, their degrees
$b_\ell$, and derived indices of AC-K25--27, define
\[
 I_{399}^{(0)}(u)=
 \sum_{1\leq e<\lfloor u^{399}\rfloor}\sum_\ell
 \frac{\log\#k_{O_\ell}}{b_\ell}
 \sum_{\substack{\eta\in\widehat{\operatorname{Gal}(H_u/K)}\\
                  \eta|_{\operatorname{Gal}(H_u/F_u)}=1}}
                  \max(0,i_{\ell,\eta,e}).                     \tag{AW.7}
\]
Then on the same $W$,
\[
               I_{399}^{(0)}(u)
                 =o\bigl(\exp(u^\varepsilon/2)\bigr).           \tag{AW.8}
\]
This is the entire actual $j=0$ sector: all its growing finite
characters, all coefficient primes, and every indicated Tate
weight are included.  It is not a bound for $j=1,2$.
\end{theorem}

\begin{proof}
\proofstep{1. Fixed input data and a common choice of $J$.}
Put $t(u)=\lfloor u^{C_t}\rfloor$.
The full proof \texttt{AdmissibleOfPageN.lean} shows that the
page and $t$ are $\texttt{admitted}(4,C_t,h,t)$.
For each $\psi$ set
\[
 M_\psi=\texttt{AbF.ofChar}(\psi),\quad
 B_\psi=\texttt{Bchi}(\psi),\quad
 A_{\psi,\ell}=\texttt{AbF.cyc}(\ell\psi.N).
\]
These are fixed data, even though the later fields depend on $u,q,d$.
The literal \texttt{Bchi} is finite, consists of odd primes, and
each $A_{\psi,\ell}$ contains $\mu_\ell$.
Choose $R_{0,\psi}$ using (AW.2), and choose
$R_{\psi,\ell}$ using \texttt{FixedIn}, for
$\ell\in B_\psi$.  Define
\[
 \begin{split}
 R_*&=\sum_{\psi\in\Psi}R_{0,\psi}
            +\sum_{\psi\in\Psi}\sum_{\ell\in B_\psi}R_{\psi,\ell},\\
 \eta_0&=\varepsilon/(2C_K),\qquad
 J=2^{R_*+2}(C_t/\eta_0)^{R_*}.
 \end{split}                                                 \tag{AW.9}
\]
In particular $J\geq4$ and $0<\eta_0\leq C_t$.
Duplicating an auxiliary field for different $\psi$ only enlarges
$R_*$ by a fixed amount; no disjointness is needed.

For each $\psi$, apply \texttt{PointHV} with
$j_\psi(\ell)=\operatorname{Nat.log}(\ell,R_{\psi,\ell})+3$
on $B_\psi$, using the source's natural-number logarithm,
which is also defined at $R=0$.
Choose arbitrary values outside $B_\psi$.
It supplies constants $C_\psi$ and two auxiliary primes.
All are fixed before $X,q,d$ are prescribed.

\proofstep{2. The exceptional sets are combined before choosing a window.}
This is the new finite-family step, compared with AC-G31.
The exact signature of \texttt{PartA.partA} permits
\emph{any positive $J$}; it does not require a separate formula
in terms of its own rank $R_{0,\psi}$.
The same is true of \texttt{PartB.partB}.
Apply every instance using the single $J$ in (AW.9), and take the
maximum of their finitely many thresholds.
These thresholds are uniform in the later pair $q,d$.
Specifically, \texttt{PartA} first fixes $M,R,J,C_t,\varepsilon$
and its threshold, and only afterwards quantifies over $X,q,d,L_0$.
Equation (AW.2) supplies the required $L_0$ for the prescribed pair.

For sufficiently large $X$, the given \texttt{WinPair} implies
every required $\texttt{PrimePair}(\psi.N,q,d)$,
\[
 \log(qd)\leq X^{\varepsilon/(4J)},\qquad
 q,d>2X^{C_t}+1,                                                 \tag{AW.10}
\]
and makes both primes larger than all fixed auxiliary primes and
all fixed prime divisors in the finitely many \texttt{Host.Big}
conditions.  These are the same estimates used in the complete
proof of \texttt{vwinAt\_of\_joint}:
write $Y_0=X^{1/J}$, use
$\log(qd)\leq Y_0^{\varepsilon/5}(1+(\log Y_0)^2)
\leq Y_0^{\varepsilon/4}$ eventually, and use the lower bound
$\log q>Y_0^{\varepsilon/5}/2$ (and the stronger lower bound
for $\log d$).  A fixed positive power of $X$ dominates $\log X$.
The exclusions of finitely many fixed primes are therefore uniform.

Part A supplies an exceptional set $\mathcal E_{0,\psi}$ of
cardinality at most $R_{0,\psi}$; Part B supplies
$\mathcal E_{\psi,\ell}$ of cardinality at most $R_{\psi,\ell}$.
Let $\mathcal E$ be their union.  Then $|\mathcal E|\leq R_*$.
These sets can depend on $X,q,d$ as permitted by the source.
The proof does not assume that they are fixed once and for all.
The complete elementary source proof \texttt{Window1.lean},
followed by \texttt{Window.lean}, gives $W$ as in (AW.4) with
\[
 \forall u\in[W,W^2],\quad
 \forall e\in\mathcal E,\quad e\leq u^{C_t}
                          \Longrightarrow e\leq u^{\eta_0}.
                                                                    \tag{AW.11}
\]
For clarity, its underlying induction has the following form.
For a set of at most $R$ points on the logarithmic axis, and
$B\geq1,U>0$, it finds
$U/(2^{R+2}B^R)\leq y$, $2y\leq U$, such that any marked
$a\leq v\in[y,2y]$ satisfies $Ba\leq v$.
If a mark occurs in $(U/(2B),U]$, remove it and apply the
induction below $U/(2B)$; otherwise take $y=U/2$.
Apply this with $B=C_t/\eta_0$, $U=\log X$, and marks
$\log e/C_t$ for positive $e\in\mathcal E$; $e=0$ is harmless.
Exponentiation proves (AW.11).

The condition for the union implies the condition for each subset.
Consequently every Part A conclusion and every Part B conclusion
holds on this one $W$, at every $u$ in its full squared window.
Taking $X_0$ large enough makes $X^{1/J}$ exceed all pointwise
thresholds.  No intersection of previously chosen windows occurs.

\proofstep{3. The actual pointwise automorphic estimate.}
Part A bounds the coefficient-two sum for each $E_{\psi,u}$ by
$\exp(u^{4\varepsilon/5})$ before multiplication by $\log2$.
Part B bounds every required layer in
$\texttt{LayerLe}(\psi,\ell,j_\psi(\ell),t(u),B)$ with
$B=\exp(u^{4\varepsilon/5})$.
The quantification over all models of the auxiliary fields in
\texttt{PartB} is exactly the one in \texttt{LayerLe}.
By (AW.10), \texttt{Host.Big} holds for the actual host
$(\psi,h(u),q,d)$ and $t(u)$.
Apply \texttt{PointHV} directly, as in the full proof
\texttt{ThmHV.lean}.  This avoids requiring that the common $J$
equal the single-character wrapper's \texttt{Jchi}.
Using \texttt{hostCond}, its conclusion is
\[
 \texttt{VVplusLe}\left(E_{\psi,u},t(u),
       (\operatorname{cond}(E_{\psi,u})t(u))^{C_\psi}
              +C_\psi\exp(u^{4\varepsilon/5})\right).          \tag{AW.12}
\]
The field is the literal \texttt{IsChiField}: its stabilizer is
the intersection defining $B_{h(u)}L_\psi$.
It is CM since it contains $\mathbb Q(\sqrt{-q})$ and is abelian
over $\mathbb Q$.

The retained \texttt{ThmBcFree} bounds the complementary
odd-prime eigenspaces.  The complete conductor and page proofs give
\[
 \log(\operatorname{cond}(E_{\psi,u})t(u))
                                  \leq2u^{\varepsilon/4},
 \qquad
 \texttt{minusBound}\leq K_\psi\exp(8u^{\varepsilon/4}).        \tag{AW.13}
\]
Here $K_\psi$ is fixed.  All relevant degrees are bounded by the
conductor, as in \texttt{ConductorBoundsOf\allowbreak IsChiField.lean}.
The complete \texttt{LengthEqAddEigen.lean} proof splits
the two conjugation eigenspaces at odd coefficient primes and
keeps the prime two separately.  Its \texttt{vvle\_of\_parts}
therefore bounds (AW.3) by
\[
 \begin{split}
 &\log2\exp(u^{4\varepsilon/5})
   +K_\psi\exp(8u^{\varepsilon/4})\\
 &\hspace{12mm}+\exp(2C_\psi u^{\varepsilon/4})
                   +C_\psi\exp(u^{4\varepsilon/5}).
 \end{split}                                                 \tag{AW.14}
\]
Take a maximum over the finite family to obtain (AW.5).
Its stronger little-o conclusion follows without changing
$\varepsilon$: for fixed $C$,
$\log C+Cu^{4\varepsilon/5}-u^\varepsilon/2\to-\infty$.
Thus the estimate supplies every fixed positive fraction of
the desired budget after increasing the initial threshold.

\proofstep{4. The noncyclic abelian field and its integral coefficients.}
Write
$B_u=\mathbb Q(\zeta_{2^{h(u)+2}})^+
\mathbb Q(\sqrt{-q})\mathbb Q(\sqrt{-d})$.
The exact AC-G30 computation gives
\[
 F_u=B_u\mathbb Q(\zeta_{84}),\qquad
 \operatorname{Gal}(F_u/B_u)\simeq C_2^3\times C_3.              \tag{AW.15}
\]
Restriction identifies its $24$ characters with the fixed family
just used.  Every $E_{\psi,u}=B_uL_\psi$ is an actual subfield
of $F_u$; its relative cyclic degree is at most six.

Fix $\ell$ and use a completion $O^{(6)}_\ell$ of
$\mathbb Z[\zeta_6]$, of degree $b^{(6)}_\ell\leq2$.
Let $\mathcal T_{\ell,\eta}(B_u,e)$ be the strict character
group over $B_u$ with this coefficient and Tate twist $-e$.
Integral Fourier, Shapiro, and the full-local/inertia restriction
squares of AC-K25 give
\[
 \left|\log\#S_\ell(F_u,e)
 -\frac1{b^{(6)}_\ell}\sum_\eta
                     \log\#\mathcal T_{\ell,\eta}(B_u,e)\right|
                                  \leq E_{F,\ell}(u,e).        \tag{AW.16}
\]
This is exact when $\ell\ne2,3$.
For each individual $\eta$, use the Fourier comparison for the
cyclic subfield cut out by $\eta$.  That comparison sums over
all characters of its cyclic quotient, one of which is $\eta$.
Nonnegativity gives
\[
 \frac1{b^{(6)}_\ell}\log\#\mathcal T_{\ell,\eta}(B_u,e)
       \leq\log\#S_\ell(E_{\eta,u},e)+E_{\eta,\ell}(u,e).
                                                                    \tag{AW.17}
\]
Again the error vanishes away from $2,3$.
This argument uses two finite Fourier comparisons; it does not
assert that a map into a larger sum of permutation modules has
finite cokernel.

At $2,3$, AC-G30-Fourier supplies the exact fixed defects for
(AW.16).  The cyclic quotient comparisons also have fixed rank,
fixed exponent, and finite action factoring through
$\mathbb Q(\zeta_{84})$.
The common fixed field $\mathbb Q(\zeta_{168})$ trivializes
these actions and every Tate twist of the finite defects.
AC-K19 applies over the totally imaginary field $B_u$.
Its bound is uniform in $e$ and is
$O(\log D_{F_u}+[F_u:\mathbb Q]+s_{F_u})$ per weight;
including the finitely many ramified rational primes gives
$s_{F_u}=O([F_u:\mathbb Q])$.
Hence the sum of all errors in (AW.16)--(AW.17), through
$e<\lfloor u^{401}\rfloor$, is
\[
 O_\varepsilon\bigl(u^{403+41\varepsilon/5}(\log u)^2\bigr).
                                                                    \tag{AW.18}
\]
This is the already proved growing-host estimate in AC-K19.
It is uniform for the prescribed pair because $u\geq X^{1/J}$.
Sum (AW.16)--(AW.17), then the $24$ bounds (AW.5).
Equation (AW.18) is little-o of the budget, proving (AW.6).

\proofstep{5. The original $j=0$ summands and coefficient ranks.}
The characters in (AW.7) are exactly those of
$\operatorname{Gal}(F_u/K)$, with exponent $3\cdot2^{h(u)}$.
Thus they are all valued in the original $O_\ell$ of AC-K25;
no new integral lattice or coefficient realization is substituted.
Apply the proof of AC-K25 directly to $F_u/K$, using those
$O_\ell$ and dividing orders by their actual
$b_\ell=\operatorname{rank}_{\mathbb Z_\ell}O_\ell$.
When $\ell\ne269$, the supplementary prime $269$ in the original
localization set can be removed for these coefficients by the
same inflation and inertia-kernel argument in AC-K25, Step 3.
When $\ell=269$, it remains a coefficient-prime place with the
full local zero condition throughout.
This gives the all-prime comparison of the sum of their strict
lengths with $\mathcal B(F_u,\lfloor u^{399}\rfloor)$,
with the proved polynomial Fourier cost at $2,3$.
The auxiliary rings $O^{(6)}_\ell$ in Step 4 are already
eliminated through the intrinsic group $S_\ell(F_u,e)$;
they never replace the original $O_\ell$ in (AW.7).

For every summand the determinant comparison (EU.2) gives
$\max(0,i_{\ell,\eta,e})\leq
\operatorname{length}_{O_\ell}H^2_{\ell,\eta}(e)$.
The absolute local comparison FL.4 may be restricted to this
subset, since its right-hand cost bounds the sum for all
characters.  Its total cost at cutoff $399$, together with the
growing Fourier cost, is
\[
 O_\varepsilon\bigl(u^{802+16\varepsilon}(\log u)^2\bigr).
                                                                    \tag{AW.19}
\]
Every coefficient prime is retained, including $2,3,269$;
the finite prime supports are those already proved in
AC-K25--27.  Equations (AW.6) and (AW.19), and monotonicity
from cutoff $401$ to $399$, prove (AW.8).
\end{proof}

\paragraph{Exact relationship with the old open questions.}
AC-QG100 records the loss in AC-G31 when one is given only
separate \texttt{VWinAt} statements.  That deduction remains
correct under its stated hypotheses.  The proof above supplies
the missing finite-family synchronization from the stronger
actual source construction, using its literal uniform
\texttt{TwoIn} and arbitrary-$J$ Part A/Part B interfaces.
It removes that source-budget loss for the abelian field $F_u$.
AC-QG101's obstruction to the stated naive change of
$\varepsilon$ is not altered: no such reparameterization is used.

\paragraph{The surviving I399 common-window condition.}
Write $I_{399}=I_{399}^{(0)}+I_{399}^{(\ne0)}$.
If a bound with fixed positive slack for $I_{399}^{(\ne0)}$
is supplied on the very same $W$ constructed above, (AW.8)
absorbs the abelian sector and gives (TC.3), after increasing
the threshold.  This is an exact conditional use of the result.
An independent assertion of the form ``there exists some
$W$ for the remainder'' does not supply this compatibility.
In particular, a squared window $[Y,Y^2]$ contained in
$[W,W^2]$, with $W>1$, forces $Y=W$: its endpoints give
$W\leq Y$ and $Y^2\leq W^2$.
There is no general operation of shrinking an already chosen
squared window to a different squared window inside it.
Nor does the source give a prescribed-window estimate in place
of its exceptional-set construction.
Accordingly (AW.8) is not asserted to weaken the existential
remainder-only form of I399, and does not close EB-Q02,
$\mathcal R_{\rm cub}$, or AM-H04.4.

\paragraph{Geometric and source scope.}
All groups are the original cohomology of arithmetic curves
with pure Tate or finite Artin--Tate coefficients.
The finite cover/Fourier dictionary is exactly AC-K19/25 and
AC-G30; no new motive, Hodge-square construction, or mixed
theta identification is claimed.  Complete applied containing
source proofs and their exact hashes are listed in the adjacent
receipt.  Deep automorphy, duality, Ferrero--Washington, and
the source's tower bound are retained at the stated interfaces;
reading their wrappers is not represented as a reproof of those
deep inputs.
